% !PRINTED 384
\weberpage{384}
\providecommand{\weberfourup}{\mathord{\tikz[baseline=-0.2ex,x=0.23em,y=0.23em,line width=0.45pt]{\draw (0,0)--(-1.05,1) (0,0)--(-0.35,1.08) (0,0)--(0.35,1.08) (0,0)--(1.05,1);}}}
\begin{center}
\S.\,102.

\gesperrt{Representation of the Group.}
\end{center}

The degree of the group $P$ can be stated at once. Since there are 288 complete seven-systems, and since the elements of such a system can be permuted in $\Pi(7)$ ways, the degree of the group is:

\[
288\,\Pi(7) = 36\,\Pi(8) = 1\,451\,520.
\]

In forming the permutations of the roots we make use of two kinds of notation. First

\weq{1}{\xi_1,\ \xi_2,\ \dots,\ \xi_7,\ \xi_{ik},}

in which $i,k$ runs from 1 to 7, and the uniform notation $[i\,k]$, where $i,k$ runs from 1 to 8, and where $\xi_1,\xi_2\dots$ are to be denoted by $[1\,8],[2\,8],\dots$.

If we now permute two figures from the sequence $1,2,\dots,7$, e.\,g. 1 with 2, then $\xi_1$ passes into $\xi_2$, and by the definition of \S.\,99 $\xi_{12}$ remains unchanged, $\xi_{13}$ passes into $\xi_{23}$, and so on.

If, then, we permute the figures 1 to 7 arbitrarily in the sequence of the roots $[i\,k]$, we obtain nothing but permutations of the group $P$.

Now we have seen in \S.\,100 that, in the arrangement into syzygetic and azygetic triples, and in consequence also into complexes and complete systems, the figure 8 appears as completely equally entitled with the remaining figures 1 to 7, and since the assignment of the root $\xi_{ik}$ to the pair $\xi_i\xi_k$ [by the formula \S.\,101, (5)] depends by \S.\,99 only on this arrangement, we can also permute the figures $1,2,\dots,7,8$ without leaving the group $P$. There is therefore contained in $P$ a divisor which is isomorphic with the symmetric group of the permutations of eight elements, which therefore has the index 36, and which we shall denote by $S$.

To determine the still missing permutations of $P$, we let a new seven-system of the type $\Delta\,\weberfourup$ take the place of the seven-system $\xi_1,\dots,\xi_7$, say as follows:

% [sic] The source prints a full stop after [6 8] in the first row.
\weq{2}{\begin{pmatrix} [1\,8], & [2\,8], & [3\,8], & [4\,8], & [5\,8], & [6\,8]. & [7\,8] \\ {}[2\,3], & [3\,1], & [1\,2], & [4\,8], & [5\,8], & [6\,8], & [7\,8] \end{pmatrix},}
%CONT>>
